\begin{agradecimentos}


Eu, Lucas de Lima Spinosa dos Santos, agradeço à minha mãe, Luciana de Lima da Silva, à minha madrasta, Marta Machado, e ao meu pai, Francisco Spinosa dos Santos, pelo apoio durante toda a realização deste trabalho. Vocês me apoiaram nos períodos de dificuldade, e sou imensamente grato por isso e pela compreensão diante das mudanças em minha rotina em casa, necessárias para que este trabalho fosse concluído. Obrigado.

À minha avó, Isami Ribeiro da Costa, pelo carinho, acolhimento e apoio não apenas ao longo desta trajetória, mas também ao longo da vida. Obrigado por fazer parte da minha vida e por me permitir fazer parte da sua. Sou profundamente grato por todos os momentos que compartilhamos, por cada conversa, cada risada e por todas as lembranças que construímos juntos. Eu te amo muito, vó.

À minha mãe de coração, Tatiana Ulhoa, pelo carinho, acolhimento, conselhos e apoio ao longo dos anos, inclusive durante a realização deste trabalho. Obrigado por me incentivar a seguir em frente, mesmo nos momentos mais difíceis, quando eu já não acreditava em mim e me faltavam forças para continuar. Obrigado por todos os momentos, pelas conversas e pelo tempo que passamos juntos. Obrigado por continuar acreditando em mim. Sinto que nunca serei capaz de retribuir tudo o que fez e continua fazendo por mim, mas quero, ao menos, orgulhá-la com minhas conquistas. Te amo mais do que pode imaginar.

Eu, Rodrigo Carvalho dos Santos, expresso minha profunda gratidão à minha mãe, Raimunda Silva, e ao meu padrasto, Francismar Menezes, que me acolheram e supriram todas as minhas necessidades durante todo o curso superior, oferecendo o suporte essencial para que eu pudesse focar nos meus objetivos acadêmicos. Agradeço também ao meu pai, Milton Santos, e à minha madrasta, Claudia Castro, por me incentivarem, aconselharem e inspirarem constantemente a nunca desistir e a sempre acreditar no meu potencial, celebrando cada etapa desta jornada comigo.

Ao meu orientador, Vandor Rissoli, meu sincero agradecimento pela parceria indispensável. Ele esteve disponível, acompanhou todo o processo de desenvolvimento, orientou os caminhos da pesquisa e ajudou a refinar o trabalho até a excelência. Um agradecimento especial à minha dupla de TCC, cujo papel foi fundamental desde a concepção da ideia do projeto que desenvolvemos. Ele assumiu a frente em diversas áreas durante a execução, e nossa parceria mútua foi o motor que nos manteve motivados a superar os desafios e concluir este trabalho com sucesso.

Por fim, agradeço ao meu melhor amigo, Lucas Alves, cuja presença foi crucial para me manter motivado ao longo de todos os nossos 12 anos de amizade. Agradeço pelas conversas maduras, pelos conselhos e reforços emocionais, pelos momentos de lazer e descontração, e, acima de tudo, pela lealdade e cumplicidade que tornaram essa caminhada muito mais leve.

% Esta seção de agradecimentos é opcional, portanto, sua inclusão 
% fica a critério do(s) autor(es), que caso deseje(em) fazê-lo deverá(ão) 
% utilizar este espaço, seguindo a formatação de \textit{espaço simples e 
% fonte padrão do texto (arial ou times, tamanho 12 sem negritos, aspas ou 
% itálico}.

% \textbf{Caso não deseje utilizar os agradecimentos, deixar toda este arquivo
% em branco}.
\end{agradecimentos}
